\subsection{The large-prime tail}\label{subsec:v5-large-sieve}

We use two exact inputs from Browning--Heath-Brown~\cite{v5bhb}.
Their Lemma~2.1 says that coprime polynomials in $\mathbb Z[X_1,\ldots,X_m]$
of bounded degree and coefficient height at most $H$ have at most
\begin{equation}\label{eq:v5-box-sieve}
 O\left(\frac{V\log(VH)}{L\log L}
           +\frac{V\log(VH)}{B_{\min}}\right)
\end{equation}
integer points in a box of side parameters $B_i\geq2$ at which a prime
greater than $L\geq2$ divides every polynomial; $V=\prod B_i$.
``Coprime'' here excludes a common constant prime factor as well as a
common nonconstant factor. Their Lemma~4.1 covers the zeros modulo a
squarefree $q$ of a rank-$m$ quadratic form by at most
$(3m)^{\omega(q)}q^{m-2}$ lattices of determinant $q^{m-1}$.
The number with largest Euclidean successive minimum exceeding $T$ is
\begin{equation}\label{eq:v5-lattice-tail}
 O_m((3m)^{\omega(q)}q^{2m-3}T^{-m}).
\end{equation}
In this statement every prime dividing $q$ must be odd and preserve the
rank of the form. We include the word ``odd'' explicitly. Rank preservation
alone is not the good-prime convention of the cited lemma.

\begin{lemma}[Large common divisors on the quadric]\label{lem:v5-large-sieve}
For fixed $\kappa\ne0$ and $\epsilon,\eta\in\{1,-1\}$, the number of
integral vectors of size $O(B)$ satisfying $Q_\kappa=0$ and
$\ell\mid G_\epsilon,H_\eta$ for some prime $\ell>L$ is
\begin{equation}\label{eq:v5-large-sieve}
 O_{\kappa,\alpha}\left(
    \frac{B^{3+\alpha}}{L\log L}+B^{11/4+\alpha}\right)
 \qquad(\alpha>0).
\end{equation}
\end{lemma}
\begin{proof}
Sign changes of $a,c$ reduce to $\epsilon=\eta=1$.
Put $x_0=a+c$, $x_1=a-c$, $x=b-d$, $z=b+d$, and
$R=\kappa w^2-xz$. Then $x_0x_1=R$.
Parity restrictions can be discarded in an upper bound.

First suppose that $\ell\mid x_0,x_1$. Then $\ell^2\mid R$.
For each fixed integer $v=O(B^2)$, the equation $R=v$ has
$O_{\kappa,\alpha}(B^{1+\alpha})$ triples: for each $w$, count the
factorizations $xz=\kappa w^2-v$ by the divisor bound; the zero right-hand
side occurs for at most two $w$ and contributes $O(B)$ each.
If $R\ne0$, it has $O_\alpha(B^\alpha)$ factorizations $x_0x_1=R$.
Summing over $R=h\ell^2\ne0$ gives
$O(B^{3+2\alpha}/(L\log L))$.
The case $R=0$, including $x_0x_1=0$, contributes $O(B^{2+\alpha})$.

In the remaining cases choose $t=x_1$ or $t=x_0$ such that $\ell\nmid2t$.
Take $t=x_1$ first. Multiplication by the invertible denominator $4t^2$
turns the two norm congruences into divisibility of
\[
 P=\kappa w^2+2t,\qquad
 K=(\kappa w^2-xz+t^2+2t)^2+t^2(x+z)^2.
\]
The identity
\[
 K-(t^2+x^2)(t^2+z^2)=P(P+2t^2-2xz)
\]
shows that it suffices to sieve the coprime pair
$P,(t^2+x^2)(t^2+z^2)$. They have no common nonconstant factor over
$\mathbb Q$, because the first depends nontrivially only on $w$, whereas
the second is independent of $w$. Their contents have no common prime,
since the second has coefficient $1$ on $x^2z^2$.
For $t=x_0$ the pair is
\[
 P'=\kappa(t+2)w^2-2xz,\qquad K.
\]
If $t\ne-2$, $P'$ is an irreducible ternary quadratic form of rank three
and $K$ has nonzero constant term $t^2(t+2)^2$; homogeneity rules out
divisibility. If $t=-2$, neither $x$ nor $z$ divides $K$.
The coefficient $-2$ of $xz$ excludes a common odd content prime.

In either chart, $t\mid R$. Put
\[
 S_R=\{p:p\mid2\kappa\},\qquad
 q=\prod_{\substack{p\mid t\ p\notin S_R}}p.
\]
The form $R$ has rank three at every prime in this product, and those primes
are odd. For a fixed $q$, the number of eligible $t=O(B)$ is
$O_{\kappa,\alpha}(B^\alpha)$ after an arbitrarily small adjustment of
$\alpha$. One way to see this is to majorize the count of integers
supported on $S_R\cup\{p:p\mid q\}$ by
$B^\alpha\prod_{p\in S_R\ \text{or}\ p\mid q}(1-p^{-\alpha})^{-1}$,
and absorb its $q^{o(1)}$ factor.

Apply \eqref{eq:v5-lattice-tail} with $m=3$ to $q\mid R$.
Each covering lattice contains $q\mathbb Z^3$, so all its successive minima
are $O(q)=O(B)$. It therefore has $O(B^3/q^2)$ points in the box.
The lattices with largest minimum greater than $T$ contribute, after the
sums over $q$ and $t$,
\[
 O_{\kappa,\alpha}(B^{5+2\alpha}T^{-3}).
\]
For the other lattices choose a basis with lengths comparable to the
successive minima. The resulting rectangular box has volume $O(B^3/q^2)$
and minimum side parameter $\gg B/T$.
The substitution into the polynomial pair is invertible over $\mathbb Q$.
For $\ell\nmid2t$ its determinant, of absolute value $q^2$, is invertible
modulo $\ell$ too. Thus no new common content prime outside $2t$ appears.
Remove common content supported on $2t$, which cannot affect the primes
being counted. Coefficient heights are powers of $B$ with fixed exponents.
Formula \eqref{eq:v5-box-sieve} bounds this lattice's contribution by
\[
 O\left(\frac{B^3\log B}{q^2L\log L}
                 +\frac{B^2T\log B}{q^2}\right).
\]
There are at most $9^{\omega(q)}q$ lattices. Summation gives
$O(B^{3+3\alpha}/(L\log L)+B^{2+3\alpha}T)$.
Choose $T=B^{3/4}$ and rename the small exponent to obtain
\eqref{eq:v5-large-sieve}.
\end{proof}

\begin{proposition}[Weighted obstruction tail]\label{prop:v5-obstruction-tail}
For fixed distinct $c_j$ there are $\tau_0,C>0$ such that, for every fixed
$0<\tau\leq\tau_0$, the logarithmically weighted number of pairs
\[
 N/2\leq u\leq N,\qquad 1\leq w\leq\tau\sqrt N,
 \qquad p_j=u+c_jw^2\text{ all prime},
\]
having any of the odd common divisors forbidden in
Theorem~\ref{thm:v5-simultaneous-primes} greater than $L$, is at most
\begin{equation}\label{eq:v5-obstruction-tail}
 \frac{C\tau N^{3/2}}{L\log L}
       +O_\tau(N^{3/2-1/800}).
\end{equation}
The weight is $\prod_j\log p_j$, and the bound allows every Gaussian
generator, so it also holds for any prescribed choice of generators.
\end{proposition}
\begin{proof}
Take $B=\sqrt N$ and combine the preceding bounds at $L_0=B^{1/100}$.
The medium bound loses at most $(\log B)^k$ on inserting the weight.
The large bound with sufficiently small $\alpha$ is
$O(B^{3-1/200})$ after the logarithmic weight is inserted.
At one position use the polynomial pair $a,b^2+1$ in three variables
in \eqref{eq:v5-box-sieve}.
There are only finitely many pairs of positions and signs, and a split
rational prime has eight Gaussian generators. These factors are absorbed
in $C$. Every vector arising from prime norm is primitive, since a common
integer divisor greater than one would have its square divide that prime.
Weakening the power saving to $B^{3-1/400}$ gives the displayed estimate.
\end{proof}
